% Zone „Das kennst du schon“ – aus bank/reelle-zahlen/zone.jsonl (zusammenbau v0.1)
\einheitenkopf[zone][Kennst du schon]{Das kennst du schon}
\setcounter{aufgabe}{0}
%% TODO Grundvorstellungs-Aufgabe als erste Hauptnummer der Zone (2.8) fehlt in der Bank

%% TODO Titel: Ich-kann-Satz fehlt in der Bank (Platzhalter: Kettenname) – Nr. 1
%% TODO Anweisung: ein Satz über der ersten Teilaufgabe fehlt in der Bank – Nr. 1
\begin{aufgabe}{Bruch in Dezimalzahl umwandeln (Division), abbrechende und periodische Dezimalzahlen erkennen und mit Periodenstrich schreiben, Dezimalzahl in Bruch}
%% TODO Darstellung für schwach fehlt in der Bank (reelle-zahlen-zone-f1-v1; Abschnitt „Für schwache Schüler“)
\swz{$\frac{3}{4}$ – als Dezimalzahl?}{}
%% TODO Darstellung für schwach fehlt in der Bank (reelle-zahlen-zone-f1-v3; Abschnitt „Für schwache Schüler“)
\swz{$\frac{7}{20}$ – als Dezimalzahl?}{}
\end{aufgabe}

%% TODO Titel: Ich-kann-Satz fehlt in der Bank (Platzhalter: Kettenname) – Nr. 2
%% TODO Anweisung: ein Satz über der ersten Teilaufgabe fehlt in der Bank – Nr. 2
\begin{aufgabe}{Zahlengerade mit negativen Zahlen, Brüche und Dezimalzahlen eintragen, Zahlen vergleichen und ordnen}
%% TODO Darstellung für schwach fehlt in der Bank (reelle-zahlen-zone-f2-v1; Abschnitt „Für schwache Schüler“)
\swz{$-2$ oder $-6$ – welche Zahl ist größer?}{}
%% TODO Darstellung für schwach fehlt in der Bank (reelle-zahlen-zone-f2-v3; Abschnitt „Für schwache Schüler“)
\swz{$0{,}8$; $-0{,}9$; $0{,}85$ – aufsteigend geordnet?}{}
\end{aufgabe}

%% TODO Titel: Ich-kann-Satz fehlt in der Bank (Platzhalter: Kettenname) – Nr. 3
%% TODO Anweisung: ein Satz über der ersten Teilaufgabe fehlt in der Bank – Nr. 3
\begin{aufgabe}{Quadratwurzel als Umkehrung des Quadrierens, Quadratzahlen bis zwanzig hoch zwei, Wurzel mit dem Taschenrechner und Näherungswert, Wurzel zwischen zwei Nachbar-Quadratzahlen abschätzen, keine Wurzel aus negativer Zahl}
%% TODO Darstellung für schwach fehlt in der Bank (reelle-zahlen-zone-f3-v1; Abschnitt „Für schwache Schüler“)
\swz{$\sqrt{81}$ – welche Zahl?}{}
%% TODO Darstellung für schwach fehlt in der Bank (reelle-zahlen-zone-f3-v3; Abschnitt „Für schwache Schüler“)
\swz{$\sqrt{30}$ – zwischen welchen zwei ganzen Zahlen?}{}
\end{aufgabe}

%% TODO Titel: Ich-kann-Satz fehlt in der Bank (Platzhalter: Kettenname) – Nr. 4
%% TODO Anweisung: ein Satz über der ersten Teilaufgabe fehlt in der Bank – Nr. 4
\begin{aufgabe}{Taschenrechner}
%% TODO Darstellung für schwach fehlt in der Bank (reelle-zahlen-zone-f4-v1; Abschnitt „Für schwache Schüler“)
\swz{$\sqrt{7}$ mit dem Taschenrechner – auf zwei Stellen gerundet?}{}
%% TODO Darstellung für schwach fehlt in der Bank (reelle-zahlen-zone-f4-v3; Abschnitt „Für schwache Schüler“)
\swz{$1{,}3^2$ mit dem Taschenrechner – Anzeige?}{}
\end{aufgabe}

%% TODO Titel: Ich-kann-Satz fehlt in der Bank (Platzhalter: Kettenname) – Nr. 5
%% TODO Anweisung: ein Satz über der ersten Teilaufgabe fehlt in der Bank – Nr. 5
\begin{aufgabe}{Terme zusammenfassen}
%% TODO Darstellung für schwach fehlt in der Bank (reelle-zahlen-zone-f5-v1; Abschnitt „Für schwache Schüler“)
\swz{$4a + 3a$ – zusammengefasst?}{}
%% TODO Darstellung für schwach fehlt in der Bank (reelle-zahlen-zone-f5-v3; Abschnitt „Für schwache Schüler“)
\swz{$2{,}5y + 1{,}5y - y$ – zusammengefasst?}{}
\end{aufgabe}

%% TODO Titel: Ich-kann-Satz fehlt in der Bank (Platzhalter: Kettenname) – Nr. 6
%% TODO Anweisung: ein Satz über der ersten Teilaufgabe fehlt in der Bank – Nr. 6
\begin{aufgabe}{Potenz als Malkette mit Basis und Exponent, Potenz ausrechnen (auch negative Basis), a hoch null gleich eins und negative Hochzahl als Bruch}
%% TODO Darstellung für schwach fehlt in der Bank (reelle-zahlen-zone-f6-v1; Abschnitt „Für schwache Schüler“)
\swz{$3^3$ – als Malkette und ausgerechnet?}{}
%% TODO Darstellung für schwach fehlt in der Bank (reelle-zahlen-zone-f6-v3; Abschnitt „Für schwache Schüler“)
\swz{$(-2)^3$ – ausgerechnet?}{}
\end{aufgabe}

%% TODO Titel: Ich-kann-Satz fehlt in der Bank (Platzhalter: Kettenname) – Nr. 7
%% TODO Anweisung: ein Satz über der ersten Teilaufgabe fehlt in der Bank – Nr. 7
\begin{aufgabe}{Brüche als Exponenten lesen und mit Brüchen rechnen}
%% TODO Darstellung für schwach fehlt in der Bank (reelle-zahlen-zone-f7-v1; Abschnitt „Für schwache Schüler“)
\swz{$\frac{1}{2} + \frac{1}{2}$ – ausgerechnet?}{}
%% TODO Darstellung für schwach fehlt in der Bank (reelle-zahlen-zone-f7-v3; Abschnitt „Für schwache Schüler“)
\swz{$\frac{1}{2} + \frac{1}{3}$ – ausgerechnet?}{}
\end{aufgabe}

%% TODO Titel: Ich-kann-Satz fehlt in der Bank (Platzhalter: Kettenname) – Nr. 8
%% TODO Anweisung: ein Satz über der ersten Teilaufgabe fehlt in der Bank – Nr. 8
\begin{aufgabe}{Potenz als Malkette mit Basis und Exponent, Potenz ausrechnen (auch negative Basis), a hoch null gleich eins und negative Hochzahl als Bruch – Fehler finden}
\swz[3]{Finde den Fehler und rechne richtig: \rechnung{3^{-3} &= -27}}{}
\end{aufgabe}

%% TODO Titel: Ich-kann-Satz fehlt in der Bank (Platzhalter: Kettenname) – Nr. 9
%% TODO Anweisung: ein Satz über der ersten Teilaufgabe fehlt in der Bank – Nr. 9
\begin{aufgabe}{Potenz als Malkette mit Basis und Exponent, Potenz ausrechnen (auch negative Basis), a hoch null gleich eins und negative Hochzahl als Bruch – selbst rechnen}
%% TODO Darstellung für schwach fehlt in der Bank (reelle-zahlen-zone-f6-v7; Abschnitt „Für schwache Schüler“)
\swz{$2^{-4}$ – als Bruch?}{}
\end{aufgabe}
